% !TeX program = pdflatex
% !BIB program = biber
%-----------------------------------------------------------------
%  Course report template (group work)
%  Warsaw University of Technology, Faculty of Power and
%  Aeronautical Engineering (MEiL)
%  --
%  Based on WUT-Thesis, copyleft by Artur M. Brodzki & Piotr Woźniak
%  https://github.com/ArturB/WUT-Thesis
%  Modified by Arseni Lysak, 2026. Licensed under GPL-3.0.
%-----------------------------------------------------------------

\documentclass[
    bindingoffset=5mm,  % Binding offset
    footnoteindent=3mm, % Footnote indent
    hyphenation=true    % Hyphenation turn on/off
]{src/wut-thesis}

\graphicspath{{img/}}            % Folder with figures
\addbibresource{references.bib} % Bibliography file

%-----------------------------------------------------------------
% Extra packages. The class already loads amsmath, amssymb,
% graphicx, listings, multirow, url, hyperref and biblatex.
%-----------------------------------------------------------------
\usepackage[T1]{fontenc}
\usepackage{ellipsis}
\usepackage[scaled=0.85]{beramono}  % Monospace font for code
\usepackage{xcolor}
\usepackage{subcaption}             % Sub-figures
\usepackage{float}                  % [H] float placement
\usepackage{booktabs}               % \toprule, \midrule, \bottomrule
\usepackage{siunitx}                % Units: \SI{10}{\meter\per\second}
\usepackage{pdflscape}              % Landscape pages
\usepackage[section]{placeins}      % Keep floats inside their section
\usepackage[numbered,framed]{matlab-prettifier} % MATLAB listings
\lstset{
	inputencoding=utf8,
	extendedchars=true,
	literate={θ}{{$\theta$}}1       % Allow θ in MATLAB code
}
% Style for other code: \begin{lstlisting}[style=code, language=Python]
\definecolor{codegreen}{rgb}{0,0.6,0}
\definecolor{codegray}{rgb}{0.5,0.5,0.5}
\definecolor{codepurple}{rgb}{0.58,0,0.82}
\lstdefinestyle{code}{
	basicstyle=\ttfamily\footnotesize,
	keywordstyle=\color{blue},
	commentstyle=\color{codegreen},
	stringstyle=\color{codepurple},
	numbers=left,
	numberstyle=\tiny\color{codegray},
	breaklines=true,
	showstringspaces=false,
	tabsize=4
}
% Optional: nicer highlighting with minted, needs Python (see README)
% \usepackage{minted}
% \setminted{breaklines=true}

\raggedbottom
\emergencystretch=3em
\setcounter{tocdepth}{2}

%-----------------------------------------------------------------
% Title page
%-----------------------------------------------------------------
\facultymeil % \facultymeil or \facultyeiti
\langeng     % \langeng or \langpol

\institute{Aeronautics and Applied Mechanics} % Printed as "Institute of ..."
\worktype{Project Report}                    % Type of work, large and bold
\fieldofstudy{Course name}                   % Free-text line under the type of work
\title{Title of the Report}
\author{                                     % One line per group member
	Name Surname 123456 \\
	Name Surname 123457 \\
	Name Surname 123458
}
\supervisor{Supervisor's name}               % Printed under "Project supervisor:"
\date{\the\year}

\begin{document}
	\maketitle
	\newpage

	\tableofcontents
	\newpage

	\section{Introduction}

	What is modelled or measured, and why.
	Cite sources with \verb|\cite|, e.g.~\cite{lamport1994latex, wutthesis}.

	\section{Assumptions}

	\begin{itemize}
		\item First simplifying assumption.
		\item Second simplifying assumption.
	\end{itemize}

	\section{Model description}

	Parameters are easiest to read in a table, see Table~\ref{tab:parameters}.

	\begin{table}[htbp]
		\centering
		\caption{Example parameters.}
		\label{tab:parameters}
		\begin{tabular}{l c S[table-format=2.2] l}
			\toprule
			\textbf{Parameter} & \textbf{Symbol} & {\textbf{Value}} & \textbf{Unit} \\
			\midrule
			Length   & $L$ & 18.90 & \si{\meter} \\
			Wingspan & $b$ & 13.56 & \si{\meter} \\
			\bottomrule
		\end{tabular}
	\end{table}

	MATLAB code is typeset with \texttt{matlab-prettifier}, as in Listing~\ref{lst:matlab}:
	\begin{lstlisting}[style=Matlab-editor, caption={Example MATLAB script}, label={lst:matlab}]
t = linspace(0, 10, 1000);   % time [s]
x = exp(-0.5*t) .* sin(2*t); % damped response
plot(t, x), grid on
	\end{lstlisting}

	\section{Results}

	Figure~\ref{fig:example} shows how to include a figure.
	Put your images in the \texttt{img} folder.

	\begin{figure}[htbp]
		\centering
		\includegraphics[width=0.6\textwidth]{example-image}
		\caption{Example figure.}
		\label{fig:example}
	\end{figure}

	\section{Conclusions}

	Summary of the results.

	\printbibliography

\end{document}
